% GDZ PPN236018647,native580middle;HKUST581.
\NumberedItem{110}
\indent Theoriam nostram iam ad transcendentes ellipticas immediate applicavimus.
\par\noindent\hfill1800 Iunio 5.

\medskip\begingroup\small
\indent\textls[120]{Transcendentes ellipticae} sind (vergl. den oben bei der Nr. 105 wiedergegebenen Zettel) Integrale erster Gattung. Man wird also diese Notiz auf den Algorithmus beziehen, den \textsc{Gauss} z. B. auf den oben S. 227, 228, artt. [10.] und [12.] abgedruckten Zetteln und im art. 18. der \textit{Determinatio attractionis}, Werke III, S. 354 entwickelt hat, und der nach dem geeignet normierten Integral erster Gattung mit veränderlicher oberer Grenze als Grenzwert hinstrebt. Vergl. auch den »bilinearen Algorithmus« der Scheda An, oben S. 213.
\par\noindent\hfill\textsc{Klein.\quad Schlesinger.}
\par\endgroup

\NumberedItem{111}
\indent Rectificatio Ellipseos tribus modis diversis absoluta.
\par\noindent\hfill[1800] Iun. 10.

\medskip\begingroup\small
\indent Zwei dieser Methoden sind auf dem oben S. 227 ff. abgedruckten Zettel und zwar in den artt. [12.] und [13.] wiedergegeben. Vergl. auch die Anzeige der Abhandlung \textit{Determinatio attractionis}, Werke III, S. 360 und den art. 17. dieser Abhandlung selbst, ebenda S. 354.
\par\noindent\hfill\textsc{Klein.\quad Schlesinger.}
\par\endgroup
